% Corrected French for printed p.194; departures are append-only below.
\clearpage
\SourcePage{194}

\LayerTarget{tohoku:setup:g-module-sheaf-notation}
$H^p(G,A)$. Ainsi $\Gamma^G$ est le foncteur $h_k$ avec les notations
précédentes, on a

\LayerTarget{tohoku:definition:g-invariant-sections-sheaf}
$\Gamma^G(A)=\Gamma(A)^G$ (groupe des éléments de $\Gamma(A)$ invariants sous
$G$) $=\Gamma(A^G)$, où $A^G$ désigne le sous-faisceau de $A$ formé des germes
de sections invariantes sous $G$. (Si $G$ admet un système \textit{fini} de
générateurs, alors $A^G$ est aussi l'ensemble des éléments de $A$ invariants
sous $G$).

\LayerTarget{tohoku:definition:sheaf-group-cohomology}
Les $\mathbf H^p(G,A)$ ($-\infty<p<+\infty$) sont des faisceaux foncteurs de
$A$, qui forment le foncteur cohomologique dérivé de
$\boldsymbol{\Gamma}^G(A)=A^G$. On a des homomorphismes canoniques

\LayerTarget{tohoku:map:group-cohomology-stalk}
\[
\mathbf H^p(G,A)(x)\longrightarrow H^p(G,A(x))
\]

\LayerTarget{tohoku:remark:locality-group-cohomology}
qui sont bijectifs dans tous les cas que nous aurons à considérer, et en tous
cas si $G$ est fini (en vertu de th. 4.2.2 ou du corollaire 1 à la prop. 4.2.4,
au choix!). Ces foncteurs sont de nature locale, c'est-à-dire permutent à
l'opération de restriction à un ouvert.

\LayerTarget{tohoku:definition:equivariant-cohomology}
Les $H^p(X;G,A)$ forment le foncteur cohomologique dérivé de $\Gamma^G(A)$.
On a le

\LayerTarget{tohoku:theorem:4.4.1}
\noindent\textsc{Théorème 4.4.1.} \textit{Il existe deux foncteurs spectraux
cohomologiques sur la catégorie des $G$-faisceaux, aboutissant au foncteur
gradué $(H^m(X;G,A))$, et dont les termes initiaux sont respectivement}

\LayerTarget{tohoku:eq:4.4.1}
\[
\tag{4.4.1}
\begin{aligned}
\mathrm I_2^{p,q}(A)&=H^p(X,\mathbf H^q(G,A)),\\
\mathrm {II}_2^{p,q}&=H^p(G,H^q(X,A)).
\end{aligned}
\]

\LayerTarget{tohoku:construction:operator-complex-spectral-sequences}
\textit{Ce sont aussi les deux suites spectrales du ``complexe à opérateurs''
$C(A)$, où $C(A)$ est la résolution canonique de $A$, ou les deux suites
spectrales de hypercohomologie de l'espace $X$ par rapport au complexe
$C(G,A)$ des ``cochaines de $G$ à coëfficients dans le faisceau $A$''. Ces
deux suites spectrales sont triviales si $G$ opère trivialement sur $A$ et si
de plus $A$ est un espace vectoriel sur le corps $k$; alors on a un
isomorphisme canonique :}

\LayerTarget{tohoku:eq:4.4.2}
\[
\tag{4.4.2}
H^n(X;G,A)=
\bigoplus_{p+q=n}
\mathrm{Hom}_k(H_p(G,k),H^q(X,A)).
\]

\textit{Plus généralement, supposons seulement que $G$ opère trivialement sur
$A$, alors le corollaire 2 à la prop. 4.3.2 donne une suite exacte canonique :}

\LayerTarget{tohoku:eq:4.4.3}
\[
\tag{4.4.3}
\begin{aligned}
0\longrightarrow{}&
\bigoplus_{p+q=n-1}
\mathrm{Ext}_{\mathbf Z}^1(H_p(G,\mathbf Z),H^q(X,A))\\
\longrightarrow{}&H^n(X;G,A)
\longrightarrow
\bigoplus_{p+q=n}
\mathrm{Hom}_{\mathbf Z}(H_p(G,\mathbf Z),H^q(X,A))
\longrightarrow0.
\end{aligned}
\]

\LayerTarget{tohoku:remark:4.4.1:1}
\noindent\textsc{Remarques.} 1. Dans cette dernière formule, nous avions
sous-entendu que $k=\mathbf Z$. Cependant cela n'a pas d'importance, car si
$k$ est quelconque et $A$ est un faisceau de $k$-modules où $G$ opère, alors
les $H^p(G,A)$ et $H^p(X;G,A)$ sont les mêmes, qu'on regarde $A$ comme un
faisceau de $k(G)$-modules ou comme un faisceau de $\mathbf Z(G)$-modules.

\LayerTarget{tohoku:proof:remark-4.4.1:1}
Pour le voir on est ramené à prouver que si $A$ est un faisceau injectif de
$k(G)$-modules, alors
$\mathrm{Ext}_{\mathbf Z(G)}^n(X;\mathbf Z,A)$ et
$\mathbf{Ext}_{\mathbf Z(G)}^n(\mathbf Z,A)$ (où nous écrivons $\mathbf Z$ et
$\mathbf Z(G)$ pour les faisceaux constants qu'ils définissent) sont nuls pour
$n>0$. Pour le premier, cela résulte de la suite spectrale $II$, qui montre
qu'il est isomorphe à $H^n(G,\Gamma(A))$, qui est nul puisque $\Gamma(A)$ est
un $k$-module injectif (pour lequel le résultat invoqué est bien connu).

\TohokuContinuation{195}
